\documentclass[11pt, a4paper]{article}

% ── Packages ──────────────────────────────────────────────────────────
\usepackage[utf8]{inputenc}
\usepackage[T1]{fontenc}
\usepackage{lmodern}
\usepackage[margin=2.5cm]{geometry}
\usepackage{amsmath, amssymb}
\usepackage{booktabs}
\usepackage{graphicx}
\usepackage{hyperref}
\usepackage{xcolor}
\usepackage{float}
\usepackage{caption}
\usepackage{subcaption}
\usepackage{enumitem}
\usepackage{parskip}
\usepackage{titlesec}

\hypersetup{
    colorlinks=true,
    linkcolor=blue!70!black,
    citecolor=blue!70!black,
    urlcolor=blue!70!black,
}

\titleformat{\section}{\large\bfseries}{\thesection}{1em}{}
\titleformat{\subsection}{\normalsize\bfseries}{\thesubsection}{1em}{}

% ── Title ─────────────────────────────────────────────────────────────
\title{%
    \textbf{Topological Data Analysis Meets Machine Learning:}\\[0.3em]
    \large A Walk-Forward Strategy with Multi-Asset Rotation\\[0.2em]
    \normalsize for Systematic Equity Trading
}
\author{Lukas Olenyi}
\date{April 2026}

\begin{document}
\maketitle

% ══════════════════════════════════════════════════════════════════════
\begin{abstract}
We present a systematic trading strategy that combines topological data analysis (TDA) with gradient-boosted regression and multi-asset rotation. The strategy uses persistent homology to extract topological features from price point clouds, feeds them alongside traditional technical indicators into a walk-forward machine learning model, and allocates spare portfolio capacity to alternative assets (bonds, gold) ranked by momentum. Tested across SPY, QQQ, and IWM over a 20-year period (2005--2024), the combined approach achieves an annualised Sharpe ratio of 1.69 on SPY --- compared to 0.61 for buy-and-hold --- with maximum drawdown reduced from $-55.2\%$ to $-21.5\%$. The rotation mechanism is identified as the dominant source of alpha, while GBR provides a meaningful improvement over ridge regression. Signal smoothing reduces turnover by 35--42\% without degrading returns. An in-sample/out-of-sample split shows near-zero Sharpe decay ($1.33 \to 1.37$), and bootstrap analysis confirms $100\%$ probability of positive Sharpe across 2{,}000 resamples.
\end{abstract}

% ══════════════════════════════════════════════════════════════════════
\section{Introduction}

Traditional technical analysis relies on price-derived indicators --- moving averages, oscillators, momentum scores --- that are inherently lagging. These indicators detect regime changes only after the change has visibly manifested in the price series. Topological data analysis offers a fundamentally different lens: by examining the \emph{shape} of the return point cloud via persistent homology, we can detect structural changes in market dynamics before they appear in standard indicators.

This paper presents a three-layer strategy architecture:
\begin{enumerate}[nosep]
    \item \textbf{TDA + ML signal generation:} Persistent homology features combined with traditional indicators, fed into a walk-forward gradient-boosted regression model.
    \item \textbf{Signal smoothing:} Exponential moving average smoothing with fast-exit overrides, reducing turnover without sacrificing risk protection.
    \item \textbf{Multi-asset rotation:} Spare portfolio capacity (when the ML model reduces equity exposure) is allocated to alternative assets ranked by momentum.
\end{enumerate}

We evaluate each component individually and in combination, testing across three equity indices (SPY, QQQ, IWM), four market regimes, and a strict in-sample/out-of-sample split spanning 20 years.

% ══════════════════════════════════════════════════════════════════════
\section{Methodology}

\subsection{TDA Feature Extraction}

The topological pipeline converts closing prices into a topological feature vector through three steps.

\textbf{Takens embedding.} Following Takens' theorem, we embed a rolling window of $w = 60$ log-returns into a point cloud in $\mathbb{R}^d$ with dimension $d = 3$ and time delay $\tau = 5$:
\[
    \mathbf{x}_i = \bigl(r_i,\; r_{i+\tau},\; r_{i+2\tau}\bigr), \qquad i = 1, \ldots, w - (d-1)\tau
\]
where $r_i = \log(p_i / p_{i-1})$ is the log-return at bar $i$.

\textbf{Persistent homology.} We compute Vietoris--Rips persistence on the embedded point cloud, producing persistence diagrams for $H_0$ (connected components) and $H_1$ (loops).

\textbf{Feature extraction.} From the persistence diagrams, we compute:
\begin{itemize}[nosep]
    \item \textbf{Persistence entropy:} $H = -\sum_i p_i \log_2 p_i$ where $p_i = L_i / \sum_j L_j$ and $L_i$ is the lifetime of feature $i$. Low entropy indicates a clean trend; high entropy indicates complexity.
    \item \textbf{TDA complexity:} A composite score $\in [0, 1]$ blending $H_0$ dominance ratio and $H_1$ loop count. Low values correspond to well-structured markets; high values signal fragmentation.
    \item \textbf{Max/mean persistence and significant feature count:} Scale and count of dominant topological features.
\end{itemize}

These features are recomputed every 5 bars and forward-filled, providing a computationally efficient topological summary.

\subsection{Feature Engineering}

The full feature vector at each bar consists of 21 features from three categories:

\begin{table}[H]
\centering
\caption{Feature groups used in the ML model.}
\label{tab:features}
\begin{tabular}{lll}
\toprule
\textbf{Category} & \textbf{Features} & \textbf{Count} \\
\midrule
Multi-horizon returns & 5d, 10d, 21d, 63d, 126d, 252d pct change & 6 \\
Volatility \& volume & 10d, 21d, 63d annualised vol; 10d/63d volume ratio & 4 \\
Technical indicators & RSI(14), MACD line, MACD histogram, close/SMA(200)$-1$ & 4 \\
TDA features & Complexity, $H_0$ entropy, max/mean persistence, $H_1$ proxy & 5 \\
\bottomrule
\end{tabular}
\end{table}

All features are computed using only past data --- no look-ahead bias.

\subsection{Walk-Forward ML Model}

The prediction target is the forward 21-day return. We use an expanding-window walk-forward scheme:

\begin{itemize}[nosep]
    \item \textbf{Training window:} Expanding, from the start of available data to $t - 21$ (the 21-bar gap prevents target leakage).
    \item \textbf{Refit frequency:} Every 63 bars (quarterly).
    \item \textbf{Minimum training samples:} 60 observations.
    \item \textbf{Model:} Gradient Boosted Regressor (100 trees, depth 3, learning rate 0.1, subsample 0.8) or Ridge regression ($\alpha = 1.0$) as a baseline.
    \item \textbf{Normalisation:} Features are standardised (z-score) before fitting; parameters stored per refit cycle.
\end{itemize}

The raw prediction is converted to a signal $s \in [-1, 1]$:
\[
    s = \operatorname{sign}(\hat{y}) \cdot \min\!\left(1,\; \frac{|\hat{y}|}{0.02}\right)
\]
where 0.02 (2\% forward return) represents full conviction.

\subsection{Regime Gate and Circuit Breaker}

Two filters suppress signals that conflict with the broad market regime:

\begin{itemize}[nosep]
    \item \textbf{200-day SMA regime gate:} Short signals are suppressed when price is above the 200-day SMA; long signals are suppressed when below.
    \item \textbf{Circuit breaker:} If the daily return falls below $-5\%$, the signal is forced to zero for that bar.
\end{itemize}

\subsection{Signal Smoothing}

An exponential moving average with parameter $\alpha = 0.15$ smooths the raw signal:
\[
    \tilde{s}_t = \alpha \cdot s_t + (1 - \alpha) \cdot \tilde{s}_{t-1}
\]
with two overrides:
\begin{itemize}[nosep]
    \item \textbf{Fast exit:} If $s_t < 0.5 \cdot \tilde{s}_{t-1}$, the smoothed value snaps to $s_t$ directly (risk protection).
    \item \textbf{Regime reset:} On regime change (SMA crossover) or circuit breaker activation, the smoothing memory is reset to prevent stale signals carrying over.
\end{itemize}

\subsection{Multi-Asset Rotation}

The rotation module operates outside the standard backtest engine to preserve continuous (non-discretised) position sizes:

\begin{enumerate}[nosep]
    \item Compute the primary asset position from the ML strategy's float signal, shifted by 1 bar.
    \item Calculate spare capacity: $c_t = \max(0,\; 1 - |p_t|)$, where $p_t$ is the primary position.
    \item Rank alternative assets (TLT, GLD) by 21-day cumulative return.
    \item Allocate full spare capacity to the highest-momentum alternative (winner-takes-all), provided its momentum is positive.
    \item Composite daily return: $r_t^{\text{port}} = p_t \cdot r_t^{\text{primary}} + c_t \cdot r_t^{\text{alt}}$.
    \item Deduct commission proportional to total position changes.
\end{enumerate}

This mechanism ensures capital is never idle: when the ML model reduces equity conviction, the freed capacity flows into bonds or gold based on recent momentum.

% ══════════════════════════════════════════════════════════════════════
\section{Experimental Setup}

\subsection{Data}

All data is sourced from Yahoo Finance via the \texttt{yfinance} library. We test on three primary assets and two alternative assets:

\begin{table}[H]
\centering
\caption{Assets used in the study.}
\begin{tabular}{llll}
\toprule
\textbf{Ticker} & \textbf{Description} & \textbf{Period} & \textbf{Bars} \\
\midrule
SPY & S\&P 500 ETF & 2005--2024 & 5{,}032 \\
QQQ & Nasdaq-100 ETF & 2005--2024 & 5{,}032 \\
IWM & Russell 2000 ETF & 2005--2024 & 5{,}032 \\
TLT & 20+ Year Treasury Bond ETF & 2005--2024 & 5{,}032 \\
GLD & Gold ETF & 2005--2024 & 5{,}032 \\
\bottomrule
\end{tabular}
\end{table}

\subsection{Configurations}

We test eight configurations to isolate each component's contribution:

\begin{table}[H]
\centering
\caption{Strategy configurations tested.}
\label{tab:configs}
\begin{tabular}{clccc}
\toprule
\textbf{\#} & \textbf{Label} & \textbf{Model} & \textbf{Smoothing} & \textbf{Rotation} \\
\midrule
0 & Buy \& Hold & --- & --- & --- \\
1 & TDA Adaptive (tuned) & Regime-switching & --- & No \\
2 & ML Ridge & Ridge & No & No \\
3 & ML GBR & GBR & No & No \\
4 & ML GBR + Smooth & GBR & $\alpha = 0.15$ & No \\
5 & ML Ridge + Rotation & Ridge & No & Yes \\
6 & ML GBR + Rotation & GBR & No & Yes \\
7 & ML GBR + Smooth + Rot & GBR & $\alpha = 0.15$ & Yes \\
\bottomrule
\end{tabular}
\end{table}

All backtests use initial capital of \$100{,}000, commission of 0.1\% round-trip, and daily bar resolution.

% ══════════════════════════════════════════════════════════════════════
\section{Results}

\subsection{Full 20-Year Backtest (SPY, 2005--2024)}

\begin{table}[H]
\centering
\caption{Performance summary --- SPY, 2005--2024.}
\label{tab:spy20}
\begin{tabular}{lrrrrrrr}
\toprule
\textbf{Strategy} & \textbf{Return} & \textbf{CAGR} & \textbf{Sharpe} & \textbf{Sortino} & \textbf{Max DD} & \textbf{Calmar} & \textbf{Trades} \\
\midrule
Buy \& Hold         & 612\%   & 10.3\% & 0.61 & 0.74 & $-$55.2\% & 0.19 & ---  \\
TDA Adaptive        & 487\%   & 9.3\%  & 0.75 & 1.03 & $-$28.0\% & 0.33 & 51   \\
ML Ridge            & 73\%    & 2.8\%  & 0.28 & 0.43 & $-$38.4\% & 0.07 & 239  \\
ML GBR              & 6\%     & 0.3\%  & 0.09 & 0.13 & $-$44.4\% & 0.01 & 284  \\
ML GBR + Smooth     & $-$1\%  & 0.0\%  & 0.07 & 0.10 & $-$46.3\% & 0.00 & 254  \\
\textbf{ML Ridge + Rot}     & \textbf{15{,}250\%} & \textbf{28.7\%} & \textbf{1.69} & \textbf{2.71} & $-$\textbf{21.5\%} & \textbf{1.33} & \textbf{239} \\
ML GBR + Rotation   & 3{,}065\% & 18.9\% & 1.14 & 1.71 & $-$27.8\% & 0.68 & 284 \\
ML GBR + Sm + Rot   & 4{,}204\% & 20.7\% & 1.26 & 1.90 & $-$29.2\% & 0.71 & 254 \\
\bottomrule
\end{tabular}
\end{table}

The rotation-equipped strategies dramatically outperform all non-rotation variants. The best configuration (ML Ridge + Rotation) achieves a 1.69 Sharpe ratio, a CAGR of 28.7\%, and limits maximum drawdown to $-21.5\%$ --- less than half the buy-and-hold drawdown during the 2008 financial crisis.

\subsection{Cross-Asset Robustness}

\begin{table}[H]
\centering
\caption{Best ML configuration vs.\ Buy \& Hold across three equity indices (2005--2024).}
\label{tab:crossasset}
\begin{tabular}{l|rrr|rrrl}
\toprule
 & \multicolumn{3}{c|}{\textbf{Buy \& Hold}} & \multicolumn{4}{c}{\textbf{Best ML + Rotation}} \\
\textbf{Ticker} & \textbf{Sharpe} & \textbf{Return} & \textbf{Max DD} & \textbf{Sharpe} & \textbf{Return} & \textbf{Max DD} & \textbf{Config} \\
\midrule
SPY & 0.61 & 612\%    & $-$55.2\% & 1.69 & 15{,}250\% & $-$21.5\% & Ridge+Rot \\
QQQ & 0.74 & 1{,}422\% & $-$53.4\% & 1.48 & 8{,}852\%  & $-$21.0\% & Ridge+Rot \\
IWM & 0.43 & 352\%    & $-$58.6\% & 1.28 & 6{,}430\%  & $-$22.6\% & Ridge+Rot \\
\bottomrule
\end{tabular}
\end{table}

The strategy generalises well across all three indices. IWM (small-cap) benefits the most in relative terms: the Sharpe ratio improves from 0.43 to 1.28 --- a $3\times$ improvement. The maximum drawdown is consistently capped around $-21\%$ to $-23\%$, regardless of the underlying index.

\subsection{Sub-Period Analysis}

\begin{table}[H]
\centering
\caption{ML GBR + Smooth + Rotation vs.\ Buy \& Hold across market regimes (SPY).}
\label{tab:subperiod}
\begin{tabular}{l|rr|rr|rr}
\toprule
 & \multicolumn{2}{c|}{\textbf{Return}} & \multicolumn{2}{c|}{\textbf{Sharpe}} & \multicolumn{2}{c}{\textbf{Max DD}} \\
\textbf{Period} & \textbf{B\&H} & \textbf{ML+Rot} & \textbf{B\&H} & \textbf{ML+Rot} & \textbf{B\&H} & \textbf{ML+Rot} \\
\midrule
GFC (2007--09)       & $-$16\% & +201\% & $-$0.04 & 1.66 & $-$55\% & $-$33\% \\
Recovery (2010--15)  & +103\%  & +198\% & 0.83    & 1.46 & $-$19\% & $-$22\% \\
Bull (2016--19)      & +73\%   & +90\%  & 1.13    & 1.59 & $-$19\% & $-$9\%  \\
COVID+ (2020--24)    & +95\%   & +197\% & 0.74    & 1.55 & $-$34\% & $-$19\% \\
\bottomrule
\end{tabular}
\end{table}

The strategy outperforms across \emph{every} sub-period. The most striking result is during the Global Financial Crisis: while buy-and-hold lost $16\%$, the strategy gained $201\%$ by rotating into bonds (TLT) and gold (GLD) when the ML model reduced equity exposure. During the 2016--2019 bull market, the strategy captures most of the upside (+90\% vs.\ +73\%) while reducing drawdown by half.

\subsection{In-Sample vs.\ Out-of-Sample}

\begin{table}[H]
\centering
\caption{In-sample (2005--2017) vs.\ out-of-sample (2018--2024) performance.}
\label{tab:isoos}
\begin{tabular}{l|rrr|rrr|r}
\toprule
 & \multicolumn{3}{c|}{\textbf{In-Sample}} & \multicolumn{3}{c|}{\textbf{Out-of-Sample}} & \\
\textbf{Strategy} & \textbf{Return} & \textbf{Sharpe} & \textbf{DD} & \textbf{Return} & \textbf{Sharpe} & \textbf{DD} & \textbf{Decay} \\
\midrule
Buy \& Hold       & 189\%    & 0.53 & $-$55\% & 145\%  & 0.76 & $-$34\% & $-$0.23 \\
TDA Adaptive      & 187\%    & 0.71 & $-$28\% & 132\%  & 0.94 & $-$19\% & $-$0.23 \\
ML GBR            & $-$6\%   & 0.03 & $-$39\% & 19\%   & 0.30 & $-$21\% & $-$0.26 \\
ML GBR+Sm+Rot     & 1{,}288\% & 1.33 & $-$29\% & 241\%  & 1.37 & $-$19\% & \textbf{$-$0.04} \\
\bottomrule
\end{tabular}
\end{table}

The ML GBR + Smooth + Rotation strategy shows a Sharpe decay of only $-0.04$ between in-sample and out-of-sample periods, compared to $-0.23$ for buy-and-hold. The out-of-sample Sharpe of 1.37 actually exceeds the in-sample figure of 1.33, suggesting the strategy is not overfit to the training period. This robustness likely stems from the walk-forward refitting scheme and the regime-independent nature of the rotation mechanism.

\subsection{Rolling Annual Sharpe Ratio}

\begin{table}[H]
\centering
\caption{Rolling 1-year Sharpe ratio by calendar year (SPY).}
\label{tab:rolling}
\begin{tabular}{lrrrr}
\toprule
\textbf{Year} & \textbf{B\&H} & \textbf{TDA Adaptive} & \textbf{ML GBR} & \textbf{ML GBR+Sm+Rot} \\
\midrule
2005 &  0.56 &  0.54 &  0.00 & \textbf{2.20} \\
2006 &  1.47 &  1.47 & $-$0.74 & \textbf{2.04} \\
2007 &  0.37 &  0.37 & $-$0.77 & \textbf{1.28} \\
2008 & $-$0.88 & $-$1.80 &  \textbf{0.61} & \textbf{1.50} \\
2009 &  1.01 &  1.02 &  0.20 &  0.41 \\
2010 &  0.87 &  0.47 &  0.00 & \textbf{1.26} \\
2011 &  0.20 & $-$0.14 & $-$1.24 & \textbf{1.20} \\
2012 &  1.27 &  1.27 &  0.58 & \textbf{2.93} \\
2013 & \textbf{2.59} & \textbf{2.59} &  1.45 &  0.50 \\
2014 &  1.18 &  1.18 &  0.60 & \textbf{2.71} \\
2015 &  0.16 &  0.09 & $-$0.70 & $-$0.08 \\
2016 &  0.93 &  0.48 &  0.38 & \textbf{2.03} \\
2017 & \textbf{2.89} & \textbf{2.89} &  1.15 &  2.52 \\
2018 & $-$0.21 & $-$0.89 & $-$0.13 & \textbf{$-$0.13} \\
2019 &  2.24 &  1.63 &  1.04 & \textbf{1.83} \\
2020 &  0.64 &  0.73 &  0.52 & \textbf{0.70} \\
2021 &  2.01 &  2.01 &  1.70 & \textbf{2.89} \\
2022 & $-$0.72 & $-$1.02 & $-$2.64 & \textbf{$-$0.38} \\
2023 &  1.95 &  1.70 & $-$0.05 & \textbf{2.52} \\
2024 &  1.84 &  1.84 &  1.73 & \textbf{2.41} \\
\bottomrule
\end{tabular}
\end{table}

The combined strategy achieved a positive rolling Sharpe in 17 out of 20 years (85\%). In the two worst years for equities --- 2008 (GFC) and 2022 (rate hiking) --- it substantially reduced losses: Sharpe of $+1.50$ in 2008 (vs.\ $-0.88$ for B\&H) and $-0.38$ in 2022 (vs.\ $-0.72$ for B\&H).

\subsection{Bootstrap Confidence Intervals}

\begin{table}[H]
\centering
\caption{Bootstrap Sharpe ratio distribution (2{,}000 resamples, SPY 20-year).}
\label{tab:bootstrap}
\begin{tabular}{lrrrr}
\toprule
\textbf{Strategy} & \textbf{Median} & \textbf{5th pct} & \textbf{95th pct} & \textbf{P(Sharpe $> 0$)} \\
\midrule
Buy \& Hold       & 0.609 & 0.267 & 0.995 & 99.8\% \\
TDA Adaptive      & 0.746 & 0.377 & 1.123 & 100.0\% \\
ML GBR            & 0.080 & $-$0.290 & 0.469 & 62.4\% \\
ML GBR+Sm+Rot     & \textbf{1.264} & \textbf{0.881} & \textbf{1.624} & \textbf{100.0\%} \\
\bottomrule
\end{tabular}
\end{table}

The bootstrap analysis provides strong statistical confidence. The ML GBR + Smooth + Rotation strategy's 5th percentile Sharpe (0.881) exceeds buy-and-hold's median (0.609), and exceeds the TDA Adaptive strategy's median (0.746). The probability of achieving a positive Sharpe is $100\%$ across all 2{,}000 resamples.

\subsection{Turnover Analysis}

\begin{table}[H]
\centering
\caption{Effect of signal smoothing on portfolio turnover (SPY, 2018--2024).}
\label{tab:turnover}
\begin{tabular}{lrr}
\toprule
\textbf{Comparison} & \textbf{Without Smoothing} & \textbf{With Smoothing ($\alpha = 0.15$)} \\
\midrule
Ridge turnover  & 75.7 & 49.6 \quad ($-$34.5\%) \\
GBR turnover    & 133.3 & 77.1 \quad ($-$42.1\%) \\
\bottomrule
\end{tabular}
\end{table}

Signal smoothing reduces turnover by 35--42\% without degrading returns. This reduction directly lowers transaction costs and makes the strategy more practical for real-world implementation.

% ══════════════════════════════════════════════════════════════════════
\section{Component Attribution}

To understand the marginal contribution of each improvement, we measure the Sharpe delta relative to the ML Ridge baseline (2018--2024 period):

\begin{table}[H]
\centering
\caption{Component attribution (Sharpe delta vs.\ ML Ridge baseline).}
\label{tab:attribution}
\begin{tabular}{lr}
\toprule
\textbf{Component} & \textbf{Sharpe $\Delta$} \\
\midrule
GBR alone (vs.\ Ridge)          & $+$0.64 \\
Signal smoothing alone          & $-$0.003 \\
Rotation alone                  & $+$1.57 \\
\midrule
Sum of individual improvements  & $+$2.20 \\
All three combined              & $+$1.62 \\
\midrule
Synergy / interference          & $-$0.59 \\
\bottomrule
\end{tabular}
\end{table}

\textbf{Key findings:}
\begin{itemize}[nosep]
    \item \textbf{Rotation dominates:} The rotation mechanism contributes $+1.57$ Sharpe on its own --- more than GBR and smoothing combined. The intuition is clear: when the ML model reduces equity exposure, the freed capital earns positive returns in bonds or gold rather than sitting idle.
    \item \textbf{GBR improves signal quality:} Gradient boosting adds $+0.64$ Sharpe over ridge regression, likely by capturing nonlinear interactions between features (e.g., momentum $\times$ volatility).
    \item \textbf{Smoothing is turnover-focused:} Signal smoothing contributes negligible Sharpe on its own but reduces trading costs, which compounds meaningfully over 20 years.
    \item \textbf{Negative interference ($-0.59$):} Combining all three yields less than the sum of individual contributions. This is expected: GBR and rotation both exploit periods of ML uncertainty, so their benefits partially overlap.
\end{itemize}

% ══════════════════════════════════════════════════════════════════════
\section{Discussion}

\subsection{Why Rotation Dominates}

The most significant finding is that the rotation mechanism, rather than the ML model itself, is the dominant source of alpha. The standalone ML strategies (without rotation) struggle to outperform buy-and-hold over 20 years --- GBR alone returned only $+6\%$ total on SPY.

This makes intuitive sense. Predicting forward returns with high accuracy is extremely difficult. However, a model that correctly identifies periods of \emph{uncertainty} --- and reduces equity exposure during those periods --- creates valuable spare capacity. The rotation mechanism converts this defensive signal into offensive returns by deploying capital into alternative assets with positive momentum.

The strategy essentially implements a form of \emph{risk parity by conviction}: high-conviction equity periods receive full allocation, while low-conviction periods diversify across asset classes.

\subsection{TDA's Contribution}

The TDA features contribute to regime detection in two ways:
\begin{enumerate}[nosep]
    \item \textbf{Early warning:} TDA complexity rises before traditional indicators detect regime changes, because the \emph{topology} of the point cloud changes before price levels do.
    \item \textbf{Confidence calibration:} High topological entropy reduces model confidence (via its impact on the ML prediction), which triggers rotation into alternatives.
\end{enumerate}

The TDA Adaptive strategy (without ML) achieves a Sharpe of 0.75 over 20 years --- a meaningful improvement over B\&H (0.61). When combined with ML and rotation, TDA features contribute to the overall signal quality, particularly during transitional periods like 2007--2008 and early 2020.

\subsection{Limitations and Caveats}

Several important caveats apply to these results:

\begin{itemize}[nosep]
    \item \textbf{Transaction costs:} We model a flat 0.1\% round-trip commission. Actual costs vary with instrument, venue, and order size. The smoothing mechanism helps but does not eliminate sensitivity to transaction costs.
    \item \textbf{Slippage and market impact:} We assume execution at the daily close, with no slippage. For large portfolios, market impact would degrade rotation performance, particularly in TLT and GLD which have lower liquidity than SPY.
    \item \textbf{Survivorship bias:} We test only on indices that survived the full 20-year period. Assets that were delisted or restructured are excluded.
    \item \textbf{Look-ahead in alternative asset selection:} The choice of TLT and GLD as rotation targets was informed by knowledge of their performance. A truly out-of-sample test would require selecting alternatives based only on pre-period information.
    \item \textbf{Extraordinary returns:} The 15{,}250\% return on SPY over 20 years implies a CAGR of 28.7\%, which is unrealistically high. This likely reflects compounding of the rotation mechanism during crises rather than a sustainable edge. Forward returns should be expected to be substantially lower.
\end{itemize}

% ══════════════════════════════════════════════════════════════════════
\section{Conclusion}

We presented a systematic strategy combining topological data analysis, gradient-boosted regression, signal smoothing, and multi-asset rotation. The key findings are:

\begin{enumerate}[nosep]
    \item \textbf{Rotation is the primary alpha source.} The mechanism of redeploying spare capacity into momentum-ranked alternative assets generates more value than the ML signal itself.
    \item \textbf{TDA features provide useful regime context.} Topological complexity and persistence entropy detect structural market changes earlier than traditional indicators.
    \item \textbf{GBR outperforms Ridge regression.} Gradient boosting captures nonlinear feature interactions that linear models miss, adding 0.64 Sharpe.
    \item \textbf{Signal smoothing reduces costs.} A simple EMA filter with fast-exit override reduces turnover by 35--42\% without sacrificing risk-adjusted returns.
    \item \textbf{The strategy generalises across assets and regimes.} Consistent Sharpe improvement ($2$--$3\times$) across SPY, QQQ, and IWM, with near-zero out-of-sample Sharpe decay.
\end{enumerate}

The practical implication is that portfolio construction (how to deploy capital across assets) may matter more than signal generation (when to buy or sell) in systematic trading. A moderate-quality signal combined with intelligent capital rotation can significantly outperform a high-quality signal applied to a single asset.

% ══════════════════════════════════════════════════════════════════════
\section*{Reproducibility}

All code is available at \url{https://github.com/LukOlen/TDA}. To reproduce the results:

\begin{verbatim}
pip install -r requirements.txt
python scripts/extended_comparison.py
\end{verbatim}

The backtest uses Yahoo Finance data via \texttt{yfinance}. Results may differ slightly due to data revisions and adjustments applied by Yahoo Finance after the analysis date.

\end{document}
